\chapter
[\CO{[ARZ]} Arzneimittel]
{\CC{[ARZ]} Arzneimittel}
\label{chp:ARZ}


\minitoc


\section{Umgang mit Arzneimittel}



\section{Kurzbeschreibung wichtiger Notfallmedikamente}


\subsection{Herzrhythmus}

\Layout{57}{}{}{}{
\begin{tabularx}{\linewidth}{llccccX}
\hline

&

&
\multicolumn{3}{c|}{\bfseries
EKG} 
&
\bfseries Akz. Bdl. 
&
\bfseries Dosierung
\\\addlinespace\hline\addlinespace

&

&
\bfseries PR
&
\bfseries QRS
&
\bfseries QT 
&

&

\\\addlinespace\hline\addlinespace
\TabSub{Adenosin} 
&

&

&

&

&

&
Bolus: NW + KI: 

\\\addlinespace\hline\addlinespace
\bfseries Amiodaron III
&
~
&

&

&

&

&
Dos: 

\\\addlinespace\hline\addlinespace
\bfseries Digoxin 
&
~
&
${\uparrow}$
&
-
&
${\downarrow}$
&
beschl. 
&
Initial: 
0,5\,mg über 5-30~Min IV, dann 0,25\,mg alle 4-6~h (max: 1-1,5
mg am 1. Tag)

Erh:  
Je nach Kammerfrequenz und Nierenfunktion kann die IV-Gabe
wiederholt oder PO weitergefahren werden. 
\\\addlinespace\hline\addlinespace
\bfseries Flecainid 
&
~
&

&

&

&

&
Initial: 
 
Erh: 
\\\addlinespace\hline\addlinespace
\bfseries Lidocain Ib
&
~
&
-- 
&
-- 
&
-- 
&
-- 
&
\begin{compactdesc}
  \item[Dos] Pulslose KT od Asystolie: 75-100~in 1-2~min IV (1-1,5~\,\nicefrac{mg}{kg}),
  bei Bedarf alle 5~Min 50\,mg in 1~Min~IV (0,5-0,75~\,\nicefrac{mg}{kg}); max
  200-300\,mgwährend der ersten Stunde
  \item[Hämodynamisch stabile KT:]  50\,mg (0,5-0,75~\,\nicefrac{mg}{kg})~IV dann el
  CV, evt Erhaltung 1,5-5~g/24h~IV
\end{compactdesc}
\\\addlinespace\hline\addlinespace
\bfseries Procainamid Ia
&
~
&

&

&

&

&
Initial:  Erh:
\\\addlinespace\hline\addlinespace
\bfseries Propafenon Ic
&
~
&

&

&

&

&
Initial:  
\\\addlinespace\hline\addlinespace
\bfseries Propanolol II
&
~
&

&

&

&
~
&
Initial:  
\\\addlinespace\hline\addlinespace
\bfseries Sotalol III
&
~
&
${\uparrow}$
&
${\uparrow}$
&
${\uparrow}$
&
verl. 
&
 Initial:
 20\,mg über 5~min IV; bei Bedarf nach 20~min wiederholen:
20\,mg über 20~min; max 1,5~\,\nicefrac{mg}{kg}

  Erh: IV: max 1,5~\,\nicefrac{mg}{kg} IV aufgeteilt in 4 Tagesdosen\newline
PO: 160-320\,mg aufgeteilt in 2 Tagesdosen (max. 640\,mg/d)
\\\addlinespace\hline\addlinespace
\bfseries Verapamil IV
&

&
${\uparrow}$${\uparrow}$
&
--
&
--
&

&
 Initial:  
 5-10\,mg über 10~min~IV (0,075-0,15~\,\nicefrac{mg}{kg}), bei Bedarf nach
30~min wiederholen

Erh:
 IV: 0,005~\,\nicefrac{mg}{kg}/min IV kont. \newline
PO: 3x~80-160\,mg/d (max 480\,mg/d)
\\\addlinespace\hline\addlinespace
\end{tabularx}
}
{Tabelle}
{}

* Adenosin, Verapamil und Propanolol (u.a. BB) beeinflussen die akzess
Bündel nicht. Bestehen aber solche Bündel (zB WPW-Syndrom),
können diese Substanzen maligne Arrhythmien auslösen (KT, KF),
indem die AV-Überleitung verlangsamt oder blockiert wird, was den
akzess Bündel die Möglichkeit gibt, den Ventrikel direkt zu
depolarisieren (präexzitieren).


\subsubsection{Tabelle Antiarrhythmika}

\Layout{57}{}{}{}{
\begin{tabulary}{\linewidth}{LLLL}
\hline
~
 &
1. Wahl
(stabil) &
2. Wahl &
Alternativen, Bemerkungen
\\\addlinespace\hline\addlinespace
 VHFlimmern &
~
 &
~
 &
~
\\\addlinespace\hline\addlinespace
\bfseries VHFlattern &
~
 &
~
 &
~
\\\addlinespace\hline\addlinespace
\EE{Reentry-Tachy:}

\EE{AV-Knoten Reentry Tachykardie}
(160-220/min)

\EE{ Anterograde Reentry Tachykardie bei akzess.
Bündel} (150-250/min)

\EE{schmale QRS!}
&
Vagale Manöver 
&
Adenosin 
&
Gute LVEF (nach abst. Indikation):


\begin{inparaenum}
  \item Verapamil od Diltiazem
  \item BB
  \item Digoxin od Procainamid
  \item Cardioversion
\end{inparaenum}

LVEF {\textless} 40\,\% od chron HI

\begin{inparaenum}
  \item Amiodaron
  \item Digoxin
  \item Diltiazem (nicht Verapamil!)
  \item Kardioversion
\end{inparaenum}
\\\addlinespace\hline\addlinespace
\EE{Präexzitationssyndrome\textmd{ (zB WPW)}}

(= Retrograde Reentry Tachykardie bei akzess Bündel)

\EE{breites QRS!}
&
Procainamid IV

Flecainid IV 
&
Amiodaron 
&
Keine Medikamente die den AV-Knoten
verlangsamen oder blockieren:

Digoxin, Adenosin, CA (Verapamil!),
BB
\\\addlinespace\hline\addlinespace
\EE{Vorhoftachykardie (uni-/multifokal) {\textpm} Block}

(VH: 100-250/min) 
&
~
 &
~
 &
~
\\\addlinespace\hline\addlinespace
Kammertachykardie 
&
~
 &
~
 &
~
\\\addlinespace\hline\addlinespace
~
 &
~
 &
~
 &
~
\\\addlinespace\hline
\end{tabulary}
}{Tabelle Antiarrhythmika}{}

Vorhof:

AV-Knoten:

Kammer:


%

%%%%%%%%%%%%%%%%%%%%%%%%%%%%%%%%%%%%%%%%%%%%%%%%%%%%%%%%%%%
%%%%%%%%%%%%%%%%%%%%%%%%%%%%%%%%%%%%%%%%%%%%%%%%%%%%%%%%%%%
\subsubsection{Ajmalin (I)}

Gilurytmal (10\,ml\,/\,50\,mg, 2\,ml\,/\,50\,mg)

Wirkt auf VH u Ventrikel

\begin{labeling}[:]{mm}
  \item[Ind]
  \item[KI] Ventrikulare Tachykardiem Supraventrikuläre Tachykardie,
  WPW-Syndrom\footnote{Wolf-Parkinson-White-Syndrom},
  LGL-Syndrom\footnote{Lown-Gang-Levine-Syndrom}
  \item[KI]  Bradykardie, AV-Block, Schenkelblock, manifeste Herzinsuffizenz
  (ausgen. durch Arrhythmie)
  \item[Rx] Erw. 50\,mg über mind. 5 Min \Abk{IV}

  Wh.: 50\,mg nach 10 Min

  Perf.: \VonBis{0,5}{1}\MgKgH
  
  RR-Monitorkontrolle, Abbruch bei Verbreiterung der QRS-Komplexe,
  Temp.-labil
\end{labeling}

%%%%%%%%%%%%%%%%%%%%%%%%%%%%%%%%%%%%%%%%%%%%%%%%%%%%%%%%%%%
%%%%%%%%%%%%%%%%%%%%%%%%%%%%%%%%%%%%%%%%%%%%%%%%%%%%%%%%%%%
\subsubsection{Adenosin}

6\,mg\,/\,2\,ml

\begin{labeling}[:]{mm}
  \item[Ind] Diag. \& Therapie tachykarder Rhythmusstörungen:

  \EE{Reentrytachykardie} aus dem Bereich des des AV-Knotens wird
  terminiert

  \EE{SVT} wird verlangsamt

  \EE{VT}: kein Effekt
  \item[KI] \VonBis{2}{3}{\textdegree} AV-Block, Sick-Sinus-Syndrome

  \EE{CAVE}: bei Schwangerschaft, VH-Flattern/Flimmern, Asthma bronschiale
  \item[Rx] 6\,mg rasch \Abk{IV}, bei 0 Erfolg Wiederholung 12\,mg, dann 18\,mg
  nach jeweils 60 Sek.

  oder: \textbf{(3)-6-9-12} 
  
  \E{Kinder}: \VonBis{0,05}{0,15}\,\nicefrac{mg}{kg} in Schritten zu
  0,05\,\nicefrac{mg}{kg} in 2-Min-Abständen
  
  Dosisreduktion bei Z.\,n. Herz-TX (1/3-1/5) und bei \ce{Ca}-Antagonisten vom
  Nifedipin-Typ (1/4)

  Defi-Bereitschaft. Extrem kurze HWZ ({\textless}\,10 Sek). Es tritt
  ein \EE{AVB 3{\textdegree}} ein (kurz 0-Linie, keep cool)
\end{labeling}

\subparagraph{Antidot}: Theophyllin





%%%%%%%%%%%%%%%%%%%%%%%%%%%%%%%%%%%%%%%%%%%%%%%%%%%%%%%%%%%
%%%%%%%%%%%%%%%%%%%%%%%%%%%%%%%%%%%%%%%%%%%%%%%%%%%%%%%%%%%
\subsubsection{Amiodaron (\textrm{III})}

Sedacoron (1\,A\,/\,150\,mg)
\begin{labeling}[:]{mm}
\item[KI]  bei Ventrikulärer Tachykardie keine KI, sonst:

Iodallergie, Sinusbrady, Hyperthyreose, K\,↓, Torsade de pointes,
Gravid., Stillen

\item[Ind] Reanimation bei KF/pulsl. VT:

\item[Rx]  Bolus 300\,mg vor 4. Schock, überlege 1x Rep
150\,mg

\item[Ind] Therapie-resistente tachykarde supraventrikuläre oder ventrikuläre
Rhythmusstörung

\item[Rx] \,\nicefrac{mg}{kg} langsam über mind. 3 Min unter
Monitoring. Wiederholung frühestens nach 15 Min.\newline
Einmalige Infusion: 300\,mg in 250\,ml G5\,\% innerh. v. \VonBis{0,5}{2} Std

\item[Ind] WPW

\item[Rx] \VonBis{0,4}{0,8}\MgKgH

ARREST- und ALIVE-Studie beweisen besseres klinisches Outcome vs. Placebo und
Lidocain

\Ca{Cave: \textbf{LQTS}, Bradykardie, Blöcke bis Asystolie, v.\,a. in Kombination
mit anderen Antiarrhythmika (\textbf{Ca-Blocker!!!}), Steigerung ICP}
\end{labeling}


%%%%%%%%%%%%%%%%%%%%%%%%%%%%%%%%%%%%%%%%%%%%%%%%%%%%%%%%%%%
%%%%%%%%%%%%%%%%%%%%%%%%%%%%%%%%%%%%%%%%%%%%%%%%%%%%%%%%%%%
\subsubsection{Atropin}

\Z{Atropinum sulfuricum (1\,A = 0,5\,mg)}
\begin{labeling}[:]{mm}
  \item[Ind] Reanimation: PEA, Asystolie; Sinusbrady, Prophylaxe und Therapie.
  Vergiftung mit Alkylphosphaten der Hypersalivation
  \item[KI] Glaukom, Fieber, Hyperthyreose, Tachyarrhythmie,
  Prostatahypertrophie, Myasthenie
  \item[Rx] \begin{itemize}
    \item \EE{Reanimation}: Erw: 3\,mg einmalig; \E{Kinder}:
    \item \EE{Hypersalivation}: \VonBis{0,01}{0,02}\,\nicefrac{mg}{kg} (od:
    0,5\,mg Bolus, bis 3\,mg)
    \item \EE{Alkylphosphat-Intox}: Initial \VonBis{2}{4}\,mg, dann 0,5\,mg alle
    5 Min.
  \end{itemize}
  \item[An] Endobronch. Anw. bei Reanimation möglich, paradoxer Effekt bei
  Unterdosierung, auch spasmolytischer und antiemetischer Effekt, Temp.-stabil
\end{labeling}

%%%%%%%%%%%%%%%%%%%%%%%%%%%%%%%%%%%%%%%%%%%%%%%%%%%%%%%%%%%
%%%%%%%%%%%%%%%%%%%%%%%%%%%%%%%%%%%%%%%%%%%%%%%%%%%%%%%%%%%
\subsubsection{Methyldigoxin}

\Z{Lanitop (1\,ml\,/\,0,2\,mg)}

\begin{labeling}[:]{mm}
  \item[Ind] VHflattern + Vhflimmern mit schneller
  Überleitung
  \item[Rx] Erw: \VonBis{0,1}{0,2}\,mg (\VonBis{\nicefrac{1}{2}}{1} Amp zu
  0,2\,mg)
  \item[KI]  Flykosidintox, AVB \VonBis{2}{3}{\textdegree}, Sinusbradykardie,
  Hypo-K\textsuperscript{+}, Hyper-Ca\textsuperscript{++}, vor Kardioversion
  \item[Cave] \Cave{Nicht bei bereits digitalisierten Patienten, nicht nach iv-Gabe von
  Calcium, nur wenn hämodynamisch wirksam.}
\end{labeling}



%%%%%%%%%%%%%%%%%%%%%%%%%%%%%%%%%%%%%%%%%%%%%%%%%%%%%%%%%%%
%%%%%%%%%%%%%%%%%%%%%%%%%%%%%%%%%%%%%%%%%%%%%%%%%%%%%%%%%%%
\subsubsection{Esmolol (\textrm{II})}

\Z{Brevibloc (10\,ml\,/\,100\,mg, 10\,ml\,/\,2500\,mg)}

\begin{labeling}[:]{mm}
  \item[Ind]Tachyarrhythmie, insbes. SVT,
  Thyreotoxikose!, Flimmerprophylaxe nach MCI
  \item[KI]  Bradykardie, NYHA \VonBis{III}{IV}, (kardiogener) Schock, AVB
  \VonBis{2}{3}{\textdegree}, Sinusknotensyndrom, akuter Asthmaanfall
  \item[Rx] Bolus: 0,5\,\nicefrac{mg}{kg}, Rep. idem

  \E{Perfusor}: 3-18\MgKgH
  
  \EE{Monitor/RR-Kontrolle, langsam}
  
  CAVE: \textbf{keinesfalls mit Ca-Antagonisten (Verapamil) kombinieren}!
\end{labeling}

%%%%%%%%%%%%%%%%%%%%%%%%%%%%%%%%%%%%%%%%%%%%%%%%%%%%%%%%%%%
%%%%%%%%%%%%%%%%%%%%%%%%%%%%%%%%%%%%%%%%%%%%%%%%%%%%%%%%%%%
\subsubsection{Lidocain (2\,\%, \,\%x10=mg/ml) (\textrm{I})}
\begin{labeling}[:]{mm}
  \item[Ind] ischämiebed. ventr. Tachyarrh. (VT, VES),
  Breitkomplextachykardie beim MCI
  \item[Rx]Bolus 1\,\nicefrac{mg}{kg} langsam iv, dann Perf.
  \VonBis{2}{4}\MgKgH\newline
  tief endobronch. Bolus \VonBis{2}{3}\,\nicefrac{mg}{kg}
  \item[KI] AV-Block \VonBis{2}{3}{\textdegree}, Bradykardie, dekomp.
  Herzinsuffizienz (ausgen. durch Arrhythmie)
  \Fazit{Bei Flimmern ist Amiodaron 1. Wahl. (auch bei Tachykardie)}
\end{labeling}



%%%%%%%%%%%%%%%%%%%%%%%%%%%%%%%%%%%%%%%%%%%%%%%%%%%%%%%%%%%
%%%%%%%%%%%%%%%%%%%%%%%%%%%%%%%%%%%%%%%%%%%%%%%%%%%%%%%%%%%
\subsubsection{Metoprolol}


%%%%%%%%%%%%%%%%%%%%%%%%%%%%%%%%%%%%%%%%%%%%%%%%%%%%%%%%%%%
%%%%%%%%%%%%%%%%%%%%%%%%%%%%%%%%%%%%%%%%%%%%%%%%%%%%%%%%%%%
\subsubsection{Orciprenalin}

\Z{Alupent (1\,ml\,/\,0,5mg, 10ml)}
\begin{labeling}[:]{mm}
  \item[Ind] Bradykardie bei AVB 3{\textdegree}, Betablocker-Intox
  \item[Rx] Bolus 0,02\,\nicefrac{mg}{kg} (auch \Abk{SC}, \Abk{IM})

  Mit \ce{NaCl} verdünnt ml-weise unter Monitorkontrolle oder Perfusor geben.
  \item[KI] Tachykardie, Thyreotoxikose, hypertrophe obstruktive Kardiomyopathie
\end{labeling}

%%%%%%%%%%%%%%%%%%%%%%%%%%%%%%%%%%%%%%%%%%%%%%%%%%%%%%%%%%%
%%%%%%%%%%%%%%%%%%%%%%%%%%%%%%%%%%%%%%%%%%%%%%%%%%%%%%%%%%%
\subsubsection{Verapamil (\textrm{IV})}

\Z{Isoptin (2\,ml\,/\,5\,mg, 20\,ml\,/\,50\,mg) }

Sinusknoten, AV-Überleitung

\begin{labeling}[:]{mm}
  \item[Ind]
  \item[KI] SVT,
  VHFlattern/Flimmern m schneller Überleitung
  \item[Rx] Erw: \VonBis{2,5}{5}\--\,10\,mg;
  Kind: 0,1\,\nicefrac{mg}{kg};
  \E{Perfusor}: \VonBis{0,05}{0,1} \MgKgH (zB 50\,mg/50\,ml\,/\,5)
  \item[KI]  AV-Block, Herzinsuffizenz, kardiogener Schock, Sick-Sinus-Syndrome,
  WPW-Syndrom, LGL-Syndrom, Kinder {\textless} 12\,J
  
  \Cave{Keinesfalls mit Beta-Blocker}
  
  Monitor-/RR-Kontrolle
\end{labeling}


%%%%%%%%%%%%%%%%%%%%%%%%%%%%%%%%%%%%%%%%%%%%%%%%%%%%%%%%%%%
%%%%%%%%%%%%%%%%%%%%%%%%%%%%%%%%%%%%%%%%%%%%%%%%%%%%%%%%%%%
\subsection{Katecholamine}

%%%%%%%%%%%%%%%%%%%%%%%%%%%%%%%%%%%%%%%%%%%%%%%%%%%%%%%%%%%
%%%%%%%%%%%%%%%%%%%%%%%%%%%%%%%%%%%%%%%%%%%%%%%%%%%%%%%%%%%
\subsubsection{Adrenalin}

\Layout{23}{}
{
\begin{labeling}[:]{mm}
  \item[Präp] 1\,:\,10\,000: L-Adrenalin, 1\,:\,1\,000 Suprarenin.
  \item[Ind] Reanimation (schockbarer und nicht-schockbarer Rhythmus),
  Kardiogener Schock, Anaphylaxie > 3{\textdegree},  schweres Asthma bronchiale.
  \item[KI]Im Notfall keine, ansonsten:
  Hypertonie, KHK, Glaukom, Tachyarrhythmie, Hyperthyreose,
  Phäochromocytom.
  \item[\textrecipe]
  \begin{itemize}
    \item \EE{Reanimation:} Erw: 1\,mg, \EE{\E{Kinder}: 0,01\,\nicefrac{mg}{kg}}
    (1\,ml\,/\,10\,kg (bei 1:10\,000)); Wiederholung alle 3\,min.)
    \item \EE{Kardiogener Schock:}
    \VonBis{0,025}{0,3}\,{\textmu}g\,/\,kgKG\,/\,min
    \item \EE{Anaphylaxie:} 3{\textdegree} und schweres Asthma bronchiale: Erw.
    0,1\,mg \Abk{IV} oder \VonBis{0,2}{0,3}\,mg \Abk{SC}/\Abk{IM}; \E{Kinder}:
    0,01\,\nicefrac{mg}{kg}/\Abk{SC}/\Abk{IM}; Wiederholung alle \VonBis{5}{15} Min
    \item \EE{Kruppanfall:} 1:1.000 (sic!) im Vernebler 1:2 mit \ce{NaCl}
    verdünnt, 0,1\,\nicefrac{mg}{kg} inhal.
  \end{itemize}
  \item[Anm] Bei tiefer endobronchialer Applikation \VonBis{2}{3}-fache Dosis.
  Bei gleichz. Ca-Gabe \EE{Arrhythmien}; \EE{NaBic} über gl. Zugang inakt.
  Katech.; Kein Adrenalin in Akren spritzen; sehr Temp.-labil
  ({\textgreater}\,40\,{\textcelsius}), nur farblose Lsg. verw.; Je höher die
  Dosis, desto schlechter neurologisches Outcome.
  \item[Cave] Bei Betablocker-Pat. überwiegend Alpha-Wirkung.
\end{labeling}
}
{}
{./icons/under-construction.png}
{}
{}









%%%%%%%%%%%%%%%%%%%%%%%%%%%%%%%%%%%%%%%%%%%%%%%%%%%%%%%%%%%
%%%%%%%%%%%%%%%%%%%%%%%%%%%%%%%%%%%%%%%%%%%%%%%%%%%%%%%%%%%
\subsubsection{Dobutamin}

\opt{STATUS-DRAFT}{{\FontExperimental
\begin{labeling}[:]{mm}
  \item[Beschreibung:]
  \item[Ind]
  \item[KI]
  \item[Präparate:]
  \item[Ind] Kardiogener Schock mit Rückwärtsversagen
  (Stauung, RR {\textgreater}\,90)
  \item[KI] Tachyarrhythmie, Dehydratation,
  Hypertonie, Sulfit-Allergie, (subvalvuläre Aortenstenose)
  \item[Rx] \VonBis{2}{10} (20) \MugKgMin .
  \textit{Infusion:} 250\,mg in 500\,ml RLH
  \VonBis{7}{18}\,--\,36\,\nicefrac{gtt}{min}
  
  Monitor/RR-Kontrolle, wenn möglich mit Spritzenpumpe,
  Temp.-labil
  
  \Ca{Cave: bei Beta-Blocker-Pat. überwiegend Alpha-Wirkung, Inaktivierung
  durch \textbf{NaBic}!}
\end{labeling}
}}

%%%%%%%%%%%%%%%%%%%%%%%%%%%%%%%%%%%%%%%%%%%%%%%%%%%%%%%%%%%
%%%%%%%%%%%%%%%%%%%%%%%%%%%%%%%%%%%%%%%%%%%%%%%%%%%%%%%%%%%
\subsubsection{Dopamin}

\Experimental{
\begin{labeling}[:]{mm}
  \item[Beschreibung:]
  \item[Ind]
  \item[KI]
  \item[Präparate:]

5\,ml\,/\,50\,mg, 10\,ml\,/\,200\,mg, 50\,ml\,/\,250\,mg

\item[Ind] Kardiogener Schock mit Vorwärtsversagen
(RR {\textless}\,90), Nierenversagen, septischer Schock

\item[KI] Tachyarrhythmie, Dehydratation, Hypertonie,
Sulfit-Allergie, Hyperthyreose, Phäochromocytom

\item[Rx] \VonBis{2}{10} (20) {\textmu}g\,/\,kgKG\,/\,min

Monitor/RR-Kontrolle, wenn möglich mit Spritzenpumpe,
Temp.-labil
\end{labeling}
}

%%%%%%%%%%%%%%%%%%%%%%%%%%%%%%%%%%%%%%%%%%%%%%%%%%%%%%%%%%%
%%%%%%%%%%%%%%%%%%%%%%%%%%%%%%%%%%%%%%%%%%%%%%%%%%%%%%%%%%%

\subsubsection{Noradrenalin = Norepinephrin (NOR)}
\begin{labeling}[:]{mm}
  \item[Beschreibung:]
  \item[Ind]
  \item[KI]
  \item[Präparate:]
\end{labeling}


Arterenol (1\,ml\,/\,1\,mg, 5\,ml, 25\,mg)

\begin{labeling}[:]{mm}
  \item[Ind] Zustände mit niedrigem perivaskulärem Wiederstand
  (z.\,B. septischer Schock, Überdosierung von Vasodilatatoren,
  Therapie-resistente Hypotonie)

  \item[Rx] gro{\ss}e Variablilität

  \VonBis{0,01}{\EE{0,1}}\,--\,1,0\MugKgMin

  Erw: 5\,mg\,/\,50\,ml\,/\,4,2\,ml\,/\,h

  \E{Kinder}: pro 10\,kg: 5\,mg\,/\,50\,ml\,/\,0,6\,ml\,/\,	h
 \item[KI] Hypertonie, Tachykardie, Cave bei Inhalationsnarkotika,
 Cor pulmonale, KHK, Gravidität, Hyperthyreose, Glaukom

 invasive RR-Messung, Rekap kontrollieren
 \item[NW] Nekrosegefahr wenn peripher-venös od. para,
 Zentralisation, Lungenödem, Angst, Unruhe, Tachykardie, Bradykardie
\end{labeling}
%

\subsection{Kortikoide}

%

%%%%%%%%%%%%%%%%%%%%%%%%%%%%%%%%%%%%%%%%%%%%%%%%%%%%%%%%%%%
%%%%%%%%%%%%%%%%%%%%%%%%%%%%%%%%%%%%%%%%%%%%%%%%%%%%%%%%%%%
\subsubsection{Methyl-Prednisolon}

Solu-Medrol (TSA 250, 500, 1000\,mg)
\begin{labeling}[:]{mm}
  \item[Beschreibung:]
  \item[Ind]
  \item[KI]
  \item[Präparate:]
\end{labeling}
\begin{labeling}[:]{mm}\item[Ind] \item[KI] Rückenmarkstrauma

\item[Rx] Bolus: 30\,\nicefrac{mg}{kg} in 15\,min, Erhaltung:
5,4\MgKgH über 23\,h(/47\,h) (NASCIS)

\item[KI]  I. Notfall SHT%
%ÖNK Konsensus 2005: 6. Die Behandlung mit hochdosiertem Methylprednisolon ist bei gleichzeitigem Vorliegen eines Schädel-Hirn
, sonst Überempfindlichkeit, Pilz-, Herpes-, TB-Infektion, ak.
Spychosen, Ulcera

Therapiebeginn innerhalb der ersten 8 Stunden nach Querschnittslähmung


Wirkung beim WS-Trauma umstritten (ÖNK Consensus
2005\footnote{http://www.agn.at/dokumente/solumedrol.pdf}:  präklinisch nicht
empf.), wenn, dann nur Methyl-Prednisolon!

NASCIS-Schema: Bolus (30\,\nicefrac{mg}{kg}) nach {\textless}3\,h, spätestens
{\textless} 8\,h, gefolgt von Erhaltungs-Infusion 5,4\MgKgH für
23\,h (bzw. 47h, Rx erst ${\geq}$ 3\,h nach Trauma).


\end{labeling}

%%%%%%%%%%%%%%%%%%%%%%%%%%%%%%%%%%%%%%%%%%%%%%%%%%%%%%%%%%%
%%%%%%%%%%%%%%%%%%%%%%%%%%%%%%%%%%%%%%%%%%%%%%%%%%%%%%%%%%%
\subsubsection{Prednisolon}

Solu-Dacortin, Solu-Decortin (TSA: 25, 50, 250, 1000\,mg)
\begin{labeling}[:]{mm}
  \item[Beschreibung:]
  \item[Ind]
  \item[KI]
  \item[Präparate:]


\textbf{\textit{\textcolor[rgb]{0.0,0.5019608,0.0}{In}}}\textbf{ \&
}\textbf{\textit{Rx:
}}\textcolor[rgb]{0.0,0.5019608,0.0}{Stoffwechselkoma:} \VonBis{50}{100}\,mg,
\textcolor[rgb]{0.0,0.5019608,0.0}{Status asthmaticus} \VonBis{250}{500}\,mg,
\textcolor[rgb]{0.0,0.5019608,0.0}{Anaphylaxie} \VonBis{1}{3}\,g, (tox.
\textcolor[rgb]{0.0,0.5019608,0.0}{Lungenödem, Reizgasinhal.} \VonBis{1}{3}\,g)

\item[KI]  Im Notfall keine, Überempfindlichket, Pilz- Herpes- Tb-Infektion,
ak Psychosen, Ulcera

Kortikoide haben keine Sofortwirkung, Notfallindikationen sind zT
umstritten
\end{labeling}
%

\subsection{Narkose}

\begin{labeling}[:]{mm}
  \item[Präp] ---.
  \item[Beschreibung] ---.
  \item[Ind] ---.
  \item[KI] ---.
  \item[Rx] ---.
  \item[Cave] ---.
  \item[Anm] ---.
\end{labeling}

Relaxanz


%%%%%%%%%%%%%%%%%%%%%%%%%%%%%%%%%%%%%%%%%%%%%%%%%%%%%%%%%%%
%%%%%%%%%%%%%%%%%%%%%%%%%%%%%%%%%%%%%%%%%%%%%%%%%%%%%%%%%%%
\subsubsection{Cis-Atracurium}

\begin{labeling}[:]{mm}
  \item[Präp] (2,5\,ml\,/\,5mg, 5\,ml\,/\,10mg, 10\,ml\,/\,20mg)
  \item[Beschreibung]
  \item[Ind] Relaxation, Präcurarisierung, TIVA
  \item[KI] Cave bei Myasthenia gravis, neuromusk. Erkr. und
Gravid.\newline
Langsame Appl. bei kardiovask. Risiken und allergischer Diathese
\item[Rx] Intub: 0,2\,\nicefrac{mg}{kg}; Rep: 0,03\,\nicefrac{mg}{kg}\newline
\E{Perfusor}: \VonBis{3}{2}\,--\,1\MugKgMin
\item[Anm.:]Reines Isomer des Atracuriums, kaum Histaminausschüttung.
Wirk.-Eintritt \VonBis{120}{150}\,sek, Dauer \VonBis{40}{50}\,min

RepDosis wirkt etwa 20\,min.

Lagerung \VonBis{2}{8}{\textcelsius}
\end{labeling}



%%%%%%%%%%%%%%%%%%%%%%%%%%%%%%%%%%%%%%%%%%%%%%%%%%%%%%%%%%%
%%%%%%%%%%%%%%%%%%%%%%%%%%%%%%%%%%%%%%%%%%%%%%%%%%%%%%%%%%%
\subsubsection{Rocuronium}


\begin{labeling}[:]{mm}
  \item[Präp] \Z{\T{Esmeron}{\texttrademark} (50\,mg\,/\,5\,ml,
  100\,mg\,/\,10\,ml)}
  \item[Beschreibung] ---.
  \item[Ind] Relaxation, Präcurarisierung, TIVA.
  \item[KI] Myasthenia gravis, Lambert Eaton Sy., Zn.
  Poliomyelitis, strenge Indikationsstellung bei Grav./Stillen.
  \item[Rx] Intubation: \VonBis{0,3}{0,6}\,\nicefrac{mg}{kg}; Rep.:
  0,15\,\nicefrac{mg}{kg}\newline Perf.: 0,3-0,6\,\nicefrac{mg}{kg}\,/\,h.
  \item[Anm.] Bei \textit{0,6\,\nicefrac{mg}{kg}} gute Intubationsbedingungen schon nach
  1\,min, Wirkdauer ca. \VonBis{40}{50}{\textquotesingle}. Wiederholung wirkt
  etwa \VonBis{13}{20}\,min.
\end{labeling}



%%%%%%%%%%%%%%%%%%%%%%%%%%%%%%%%%%%%%%%%%%%%%%%%%%%%%%%%%%%
%%%%%%%%%%%%%%%%%%%%%%%%%%%%%%%%%%%%%%%%%%%%%%%%%%%%%%%%%%%
\subsubsection{Succinylcholin}
\begin{labeling}[:]{mm}
  \item[Präp] Lysthenon (5\,ml\,/\,100mg (2\,\%), TSA: 500\,mg + 25\,ml NaCl~0,9\,\% = 2\,\%)
  \item[Beschreibung]
  \item[Ind] Crash-Intubation
  \item[KI] Cholinesterasemangel, MH, Hyper-K+, neuromusk
Erkr, LuÖdem, perf Augenverl, schwerer LPS
\item[Anm.] \item[Rx] \VonBis{1}{1,5}~\,\nicefrac{mg}{kg};
Präcurarisieren!


Wirkungseintr \VonBis{45}{60}\,sek, Wirkdauer \VonBis{3}{4}\,min.
\item[Cave] Cave bei Kindern und Jugendlichen; irrev Herzstillstände in Pat mit
neuromusk Erkr
\end{labeling}






Benzo{\textquotesingle}s


%%%%%%%%%%%%%%%%%%%%%%%%%%%%%%%%%%%%%%%%%%%%%%%%%%%%%%%%%%%
%%%%%%%%%%%%%%%%%%%%%%%%%%%%%%%%%%%%%%%%%%%%%%%%%%%%%%%%%%%
\subsubsection{Diazepam}
\begin{labeling}[:]{mm}
  \item[Präp] Amp: Gewacalm (2\,ml\,/\,10mg), Gtt: Psychopax Drops,
  Rectal: 5\,mg, 10\,mg
  \item[Beschreibung]
  \item[Ind] Sedierung, Epilepsie, Eklampsie, Prämediaktion (CV), Rectal bei
  Krämpfen
  \item[KI] chron. Hyperkapnie, Myasthenia gravis, Intox mit
zentral dämpfenden Substanzen
\item[Rx i.V.]  0,2-0,5\,\nicefrac{mg}{kg}, maximal 20\,mg
\item[Cave] CAVE: Kumulation. Nicht bei Kindern {\textless}2\,a, auch
imAtemdepression
\item[Rx rect.] $\leq$ 15\,kg: 5\,mg, {\textgreater}15\,kg 10\,mg
\item[Anm.] Antidot: Flumazenil (Anexate)

\end{labeling}



%%%%%%%%%%%%%%%%%%%%%%%%%%%%%%%%%%%%%%%%%%%%%%%%%%%%%%%%%%%
%%%%%%%%%%%%%%%%%%%%%%%%%%%%%%%%%%%%%%%%%%%%%%%%%%%%%%%%%%%
\subsubsection{Flumazenil}

\begin{labeling}[:]{mm}
  \item[Präp] Anexate (5\,ml\,/\,0,5mg, 10\,ml\,/\,1mg)
  \item[Beschreibung] Benzoantagonist
  \item[Ind] ↑ ICP, Epi, Intox mit Trizyklika
  \item[KI]
  \item[Rx.] initial \VonBis{0,2}{0,3}\,mg, Wiederholung 0,2\,mg alle 60\,sek,
  Gesamtdosis: \VonBis{1}{2}\,mg. \E{Perfusor}: \VonBis{0,1}{0,4}\,mg/h;
  \E{Kinder}: \VonBis{4}{20}\,\nicefrac{{\textmu}g}{kg}
  \item[Anm.] evt. Auslösung von Entzugssymptomen, Temp.-stabil
\end{labeling}




\Layout{22}{Lorazepam}
{
\begin{labeling}[:]{mm}
  \item[Präp] Tavor, Temesta (1\,ml\,/\,4mg)
  \item[Beschreibung:]
  \item[Ind] Status epilepticus (1.W), Sedierung
  \item[KI] chronische Hyperkapnie, Myasthenia gravis, Intox. mit
  zentralal dämpfenden Substanzen
  \item[Rx.] Erw: \VonBis{2}{8}\,mg; Status: 4\,mg, Wiederholung mit 4\,mg nach
  10\,min\newline \E{Kinder}: 0,1\,\nicefrac{mg}{kg}
  \item[Anm.] Atemdepression, Antagonist: Anexate, Bei Erfolglosigkeit im Status
  Valproat
\end{labeling}
}
{}
{./icons/under-construction.png}
{}
{}






%%%%%%%%%%%%%%%%%%%%%%%%%%%%%%%%%%%%%%%%%%%%%%%%%%%%%%%%%%%
%%%%%%%%%%%%%%%%%%%%%%%%%%%%%%%%%%%%%%%%%%%%%%%%%%%%%%%%%%%
\Layout{22}{Midazolam}
{
\begin{labeling}[:]{mm}
  \item[Präp] Dormicum (1\,ml\,/\,5mg, 3\,ml\,/\,15mg, 5\,ml\,/\,5mg (!!), 10\,ml\,/\,50mg)
  \item[Beschreibung]
  \item[Ind] Sedierung, zerebraler Krampfanfall,
  Status epi, Narkoseeinleitungchron, Hyperkapnie, Myasthenia gravis,
  ImzdS
  \item[KI]
  \item[Rx] Sedierung: \VonBis{0,015}{0,03} \,\nicefrac{mg}{kg}\newline
  Krampf: \VonBis{0,05}{0,1}\,\nicefrac{mg}{kg}\newline
  Status: bis 0,2\,\nicefrac{mg}{kg}\newline
  Einleitung: \VonBis{0,1}{0,2}\,\nicefrac{mg}{kg}
  \item[Anm.] Atemdepression gut steuerbar, Wirkdauder individuell
  \item[Cave]  paradoxe Wirkung va bei älteren
\end{labeling}
}
{}
{./icons/under-construction.png}
{}
{}



\Layout{22}{}
{

}
{}
{./icons/under-construction.png}
{}
{}

%%%%%%%%%%%%%%%%%%%%%%%%%%%%%%%%%%%%%%%%%%%%%%%%%%%%%%%%%%%
%%%%%%%%%%%%%%%%%%%%%%%%%%%%%%%%%%%%%%%%%%%%%%%%%%%%%%%%%%%
\subsubsection{Opiate}


%%%%%%%%%%%%%%%%%%%%%%%%%%%%%%%%%%%%%%%%%%%%%%%%%%%%%%%%%%%
%%%%%%%%%%%%%%%%%%%%%%%%%%%%%%%%%%%%%%%%%%%%%%%%%%%%%%%%%%%
\subsubsection{Fentanyl}




\begin{labeling}[:]{mm}
  \item[Präp] (2\,ml\,/\,0,1mg=100{\textmu}g, 10\,ml\,/\,0,5mg)
  \item[Beschreibung:]
  \item[Ind] Opioidanalg., Einl. und Vert. d. Narkose
  \item[KI] Im Notfall keine, sonst: Stillen, Sgl., Asthma, akuter hepatischer
  Porphyrie, Myasthenia gravis, Geburtshilfe (Sektio)
  \item[Rx] Erw: \VonBis{0,05}{0,2}\,mg; \E{Kinder}: \VonBis{1}{2}\,(--\,5)
  {\textmu}g/kg \newline Dosisred. bei LPS, NINS, aälteren, schlehctem AZ.
  \item[Anm.] Tip: 0,1\,\,ml\,/\,5\,{\textmu}g. {\textmu}-Agonist, Atemdepress.,
  insbes. bei Komb. mit Benzos. In höheren Dosen Thoraxrigidität. Antagonist:
  Naloxon.
\end{labeling}





%%%%%%%%%%%%%%%%%%%%%%%%%%%%%%%%%%%%%%%%%%%%%%%%%%%%%%%%%%%
%%%%%%%%%%%%%%%%%%%%%%%%%%%%%%%%%%%%%%%%%%%%%%%%%%%%%%%%%%%
\subsubsection{Morphin (hydrochlorid)}

Vendal (1\,ml\,/\,10mg)

\begin{labeling}[:]{mm}
  \item[Präp] ---.
  \item[Beschreibung] ---.
  \item[Ind] Analgesie insbes. Herzinfarkt.
  \item[KI]  im Notfall: Asthma, spastische Schmerzen, Sgl, Stillen.
  \item[Rx] Erw: \VonBis{2,5}{10}\,mg iv, \E{Kinder}: 0,05\,\nicefrac{mg}{kg} iv; auch
  sc/im.
  \item[Cave] ---.
  \item[Anm] {\textmu}-Agonist, Atemdepression, parasympmimetische NW.
\end{labeling}







%%%%%%%%%%%%%%%%%%%%%%%%%%%%%%%%%%%%%%%%%%%%%%%%%%%%%%%%%%%
%%%%%%%%%%%%%%%%%%%%%%%%%%%%%%%%%%%%%%%%%%%%%%%%%%%%%%%%%%%
\subsubsection{Naloxon}



\begin{labeling}[:]{mm}
  \item[Präp]  Narcanti (1\,ml\,/\,0,4mg).
  \item[Beschreibung] Opiat-Antagonist .
  \item[Ind] ---.
  \item[KI] Überempfiundlichkeit, Atemdepression die nicht durch Opiate
ausgelöst wurde.
  \item[Rx] Erw: 0,4\,mg iv, im, sc\newline
\E{Kinder}: 0,1\,\nicefrac{mg}{kg} iv, im, sc.
  \item[Cave]  Cave bei Pat mit vorbestehender Herzerkr. .
  \item[Anm]  {\textmu}-Antagonist, daher Entzugssympt., auch bei Gebärenden
indiziert---.
\end{labeling}





%%%%%%%%%%%%%%%%%%%%%%%%%%%%%%%%%%%%%%%%%%%%%%%%%%%%%%%%%%%
%%%%%%%%%%%%%%%%%%%%%%%%%%%%%%%%%%%%%%%%%%%%%%%%%%%%%%%%%%%
\subsubsection{Piritramid}

\begin{labeling}[:]{mm}
  \item[Präp]  Dipidolor (2\,ml\,/\,15mg).
  \item[Beschreibung] Opiatanalgetikum mit langer Wirkdauer.
  \item[Ind] ---.
  \item[KI]  Sgl, ak hepat Porphyrie, MAO-Hemmer .
  \item[Rx]  Erw: \VonBis{1}{2}\,ml iv, \VonBis{2}{4}\,ml im\newline
KK: 0,1\,\nicefrac{mg}{kg}, SK:
\A{???}\VonBis{0,1}{0,2}\,\nicefrac{mg}{kg}\newline Wiederholung nach 30 Min .
  \item[Cave] ---.
  \item[Anm]  {\textmu}-Agonist, Wirkdauer ca 6\,h, Atemdepression, Dosisred.
  Bei älteren Pat, LPS, schlechtem AZ .
\end{labeling}


\subsection{Sonstige}

%%%%%%%%%%%%%%%%%%%%%%%%%%%%%%%%%%%%%%%%%%%%%%%%%%%%%%%%%%%
%%%%%%%%%%%%%%%%%%%%%%%%%%%%%%%%%%%%%%%%%%%%%%%%%%%%%%%%%%%
\subsubsection{Ketamin, S-}

\begin{labeling}[:]{mm}
  \item[Präp]  Ketanest-S (2\,ml = 50\,mg, 5\,ml = 25\,mg, 10\,ml =
  250\,mg).
  \item[Beschreibung] ---.
  \item[Ind] Analgesie, Anästhesie, Therapierefr. Bronchospasmus.
  \item[KI] erhöhter ICP\footnote{ICP: intrakranieller Druck} (relativ),
  penetrierende Augenverletzng, kard. Dekompensation, KHK.
  \item[\textrecipe :]
  Analgesie: \VonBis{0,1}{0,25}\,\nicefrac{mg}{kg} IV,
  \VonBis{0,25}{0,5}\,\nicefrac{mg}{kg} IM; 
  Anästhesie: \VonBis{1}{2,5}\,\nicefrac{mg}{kg} IV, 2,5\,-\,4\,\nicefrac{mg}{kg} IM;
  Bronchospasmus: \VonBis{1}{2,5}\,\nicefrac{mg}{kg} IV.
  \item[Cave] CAVE: Unterschiedliche Konzentrationen im Handel! .
  \item[Anm] Schluckreflexe / Spontanatmung ↑, Rep: halbe Initialdosis.
  V.a. bei Kindern mit Atropin kombinieren (Hypersalivation), Komb. mit Benzos
  empfohlen.
\end{labeling}



%%%%%%%%%%%%%%%%%%%%%%%%%%%%%%%%%%%%%%%%%%%%%%%%%%%%%%%%%%%
%%%%%%%%%%%%%%%%%%%%%%%%%%%%%%%%%%%%%%%%%%%%%%%%%%%%%%%%%%%
\begin{labeling}[:]{mm}
  \item[Präp] ---.
  \item[Beschreibung] ---.
  \item[Ind] ---.
  \item[KI] ---.
  \item[Rx] ---.
  \item[Cave] ---.
  \item[Anm] ---.
\end{labeling}
\subsubsection{Etomidat}

\begin{labeling}[:]{mm}
  \item[Präp]  Hypnomidate (10\,ml\,/\,20mg).
  \item[Beschreibung] Kurzhypnotikum zur Einleitung und Intub..
  \item[Ind] ---.
  \item[KI]  Im Notfall keine, sonst: Neugeborene und Sgl bis 6 Monate, strenge
Ind. ind Gravid .
  \item[Rx]  \VonBis{10}{20}\,--\,40\,mg (\VonBis{1}{2}\,A), max 60\,mg, \E{Kinder}:
  \VonBis{0,15}{0,30}\,\nicefrac{mg}{kg} .
  \item[Cave] ---.
  \item[Anm]  Bei rascher Injektion Masseterkrämpfe. Stillpause: 24\,h .
\end{labeling}



\subsubsection{Nalbuphin}

\begin{labeling}[:]{mm}
  \item[Präp] \Z{Nubain (1\,ml\,/\,10mg, 2\,ml\,/\,20mg)} .
  \item[Beschreibung]  Analgetikum f. mittelstarke Schmerzen .
  \item[Ind] ---.
  \item[KI]  Atemstörungen, ↑ ICP, Grav, Sulfitallergie .
  \item[Rx]  \VonBis{0,1}{0,2}\,\nicefrac{mg}{kg}, max. 0,25\,\nicefrac{mg}{kg}
  (Ceiling-Effekt) .
  \item[Cave] ---.
  \item[Anm]  {\textmu}-Agonist-Antagonist, daher Auslösung von Entzungssyndromen
bei süchtigen und deren Neugeborenen. .
\end{labeling}


\subsubsection{Propofol 1\,\%, 2\,\%}

\begin{labeling}[:]{mm}
  \item[Präp] \Z{(20\,ml\,/\,200mg, 50\,ml\,/\,500mg, 50\,ml\,/\,1000mg !)} .
  \item[Beschreibung] ---.
  \item[Ind] Einleitung.
  \item[KI] relativ: Hypovolämie, Grav, Stillen, Geburtshilfe,
  Fettstoffwechselstörung, {\textless}1 Mon (1\,\%), {\textless}3a (2\,\%).
  \item[Rx]
  \begin{inparaitem}
  \item Einleitung: \VonBis{2,5}{5}\,\nicefrac{mg}{kg}, Erw:
  \VonBis{2}{3}\,\nicefrac{mg}{kg}, Ältere; \VonBis{1}{1,5}\,\nicefrac{mg}{kg},
  (1\,\%: \VonBis{1}{5}\,ml\,/\.10\,kg).
  \item TIVA: Beginn: 12\MgKgH, Reduz:
  \VonBis{6}{10}\MgKgH (1\,\%:
  \VonBis{6}{12}\,ml\,/\,10\,kg\,/\,h)
  \item Sedierung: \VonBis{1}{4}\MgKgH (1\,\%:
  \VonBis{1}{4}\,ml/10kg/h).
  \end{inparaitem}
  \item[Cave] Hypotonie.
  \item[Anm]  Kardiodepressiv! ICP-Senkung .
\end{labeling}

%%%%%%%%%%%%%%%%%%%%%%%%%%%%%%%%%%%%%%%%%%%%%%%%%%%%%%%%%%%
%%%%%%%%%%%%%%%%%%%%%%%%%%%%%%%%%%%%%%%%%%%%%%%%%%%%%%%%%%%
\subsubsection{Thiopental }

\begin{labeling}[:]{mm}
  \item[Präp] ---.
  \item[Beschreibung] ---.
  \item[Ind] ---.
  \item[KI] ---.
  \item[Rx] ---.
  \item[Cave] ---.
  \item[Anm] ---.
\end{labeling}



{
Neostigmin nicht ZNS-gängig, Phyostigmin schon}





\subsection{Atmung}

%%%%%%%%%%%%%%%%%%%%%%%%%%%%%%%%%%%%%%%%%%%%%%%%%%%%%%%%%%%
%%%%%%%%%%%%%%%%%%%%%%%%%%%%%%%%%%%%%%%%%%%%%%%%%%%%%%%%%%%
\subsubsection{Fenoterol}

\begin{labeling}[:]{mm}
  \item[Präp] \Z{
  Berotec Dae (0,2mg/Hub); Partusisten (1\,ml\,/\,0,025mg, 1\,ml\,/\,0,5mg)}.
  \item[Beschreibung] ---.
  \item[Ind] ---.
  \item[KI]  Dae: im Notfall keine. Tokolyse IV: starke Blutung, vorz
  Plazentalösung, Tachyarrhythmie, Thyreotoxikose, hypertrophe obstruktive
  CMP---.
  \item[Rx]
  \begin{inparaitem}
  \item Asthma Bronchiale: \VonBis{1}{2} Hübe, Wiederholung nach
  \VonBis{5}{10}{\textquotesingle}
  \item Tokolyse mit Dosieraerosol: \VonBis{2}{5} Hübe
  \item Tokolyse \Abk{IV}: sehr langsamer Bolus: 25\,{\textmu}g.
  \E{Perfusor}: \VonBis{0,01}{0,05}\,{\textmu}g\,/\,kgKG\,/\,min., zB
  2\,mg\,/\,50\,ml\,/\,1-5 
  \end{inparaitem}.
  \item[Cave] ---.
  \item[Anm] DAe: Im Anfall meist unwirksam. IV: Kontrolle v Herzfrequenz v Mutter u
  Kind. Optional Hexoprenalin (Gynipral)..
\end{labeling}



%%%%%%%%%%%%%%%%%%%%%%%%%%%%%%%%%%%%%%%%%%%%%%%%%%%%%%%%%%%
%%%%%%%%%%%%%%%%%%%%%%%%%%%%%%%%%%%%%%%%%%%%%%%%%%%%%%%%%%%
\subsubsection{Terbutalin}

\begin{labeling}[:]{mm}
  \item[Präp]  \Z{Bricanyl 1\,ml\,/\,0,5mg}.
  \item[Beschreibung] ---.
  \item[Ind] ---.
  \item[KI] ---.
  \item[Rx] ---.
  \item[Cave] ---.
  \item[Anm] ---.
\end{labeling}



\begin{labeling}[:]{mm}\item[Ind] \item[KI] Asthmaanfall

\item[Rx] 0,005\,\nicefrac{mg}{kg} (1A/100kg) IV/SC

\begin{labeling}[:]{mm}\item[Ind] \item[KI]\end{labeling} Tokolyse

\item[Rx] 0,25\,mg ({\textonehalf} Amp)

\item[KI]  Cave bei Tachy, Thyreotox, hypertr obstr CMP, KHK, MCI, (subvalv)
Aortenstenose

Bei IV-Appl verdünnen, besser DT od Perf. SC unverdünnt.
Monitor-/RR-Kontrolle.
\end{labeling}

%%%%%%%%%%%%%%%%%%%%%%%%%%%%%%%%%%%%%%%%%%%%%%%%%%%%%%%%%%%
%%%%%%%%%%%%%%%%%%%%%%%%%%%%%%%%%%%%%%%%%%%%%%%%%%%%%%%%%%%
\subsubsection{Theophyllin}

\begin{labeling}[:]{mm}
  \item[Präp] ---.
  \item[Beschreibung] ---.
  \item[Ind] ---.
  \item[KI] ---.
  \item[Rx] ---.
  \item[Cave] ---.
  \item[Anm] ---.
\end{labeling}



\subsection{Antipsychotika}

%%%%%%%%%%%%%%%%%%%%%%%%%%%%%%%%%%%%%%%%%%%%%%%%%%%%%%%%%%%
%%%%%%%%%%%%%%%%%%%%%%%%%%%%%%%%%%%%%%%%%%%%%%%%%%%%%%%%%%%
\subsubsection{Haloperidol}

\begin{labeling}[:]{mm}
  \item[Präp]  \Z{Halodol (1\,ml\,/\,5mg)}.
  \item[Beschreibung] ---.
  \item[Ind] Psychosen, Alk-intox, Delirium tremens.
  \item[KI] Depression zebtralnervöser Funktionen, Kinder, Parkinson,
malignes Neuroleptika-Sy.
  \item[Rx]  \VonBis{5}{10}\,mg iv/im.
  \item[Cave] Long-QT-Syndrom.
  \item[Anm] Dyskinesien, Temp.-stabil.
\end{labeling}




 Antiemetikum
\begin{labeling}[:]{mm}
\item[Rx] \VonBis{0,5}{1,0}\,mg iv




\Cave{CAVE: }
\end{labeling}

%%%%%%%%%%%%%%%%%%%%%%%%%%%%%%%%%%%%%%%%%%%%%%%%%%%%%%%%%%%
%%%%%%%%%%%%%%%%%%%%%%%%%%%%%%%%%%%%%%%%%%%%%%%%%%%%%%%%%%%
\subsubsection{Psyquil (TM)}

\begin{labeling}[:]{mm}
  \item[Präp] ---.
  \item[Beschreibung] ---.
  \item[Ind] ---.
  \item[KI] ---.
  \item[Rx] ---.
  \item[Cave] ---.
  \item[Anm] ---.
\end{labeling}



%%%%%%%%%%%%%%%%%%%%%%%%%%%%%%%%%%%%%%%%%%%%%%%%%%%%%%%%%%%
%%%%%%%%%%%%%%%%%%%%%%%%%%%%%%%%%%%%%%%%%%%%%%%%%%%%%%%%%%%
\subsection{NSAR und Anti-''Au!''}

\subsubsection{Metamizol}
\begin{labeling}[:]{mm}
  \item[Präp] \Z{Novalgin (2\,ml\,/\,1000mg, 5\,ml\,/\,1500\,mg)} .
  \item[Beschreibung] ---.
  \item[Ind]  mittlere bis schwere Kolikartige Schmerzen, Paracetamol-resistentes
Fieber .
  \item[KI] Porphyrie, Glucose-6-Phosphat-Dehydrogenasemangel, Grav.
  \item[Rx] Erw: \VonBis{0,5}{1}\,g iv langsam!!\newline
Ki; 15\,\nicefrac{mg}{kg} iv langsam !!\newline
Rep nach 4 Std.
  \item[Cave] ---.
  \item[Anm] NW: RR-Abfall, roter Urin, verstärkte Blutungsneigung, Anaphylaxie,
Agranulacytose !!.
\end{labeling}



%%%%%%%%%%%%%%%%%%%%%%%%%%%%%%%%%%%%%%%%%%%%%%%%%%%%%%%%%%%
%%%%%%%%%%%%%%%%%%%%%%%%%%%%%%%%%%%%%%%%%%%%%%%%%%%%%%%%%%%
\subsubsection{Methylergometrin}

\begin{labeling}[:]{mm}
  \item[Präp] ---.
  \item[Beschreibung] ---.
  \item[Ind] ---.
  \item[KI] ---.
  \item[Rx] ---.
  \item[Cave] ---.
  \item[Anm] ---.
\end{labeling}


%%%%%%%%%%%%%%%%%%%%%%%%%%%%%%%%%%%%%%%%%%%%%%%%%%%%%%%%%%%
%%%%%%%%%%%%%%%%%%%%%%%%%%%%%%%%%%%%%%%%%%%%%%%%%%%%%%%%%%%
\subsubsection{Paracetamol}

\begin{labeling}[:]{mm}
  \item[Präp] ---
  \item[Beschreibung] ---
  \item[Ind]  Fieber, leichte bis mittelstarke Schmerzen .
  \item[KI] Überempf, G-6-P-DHG-Mangel, bei schwerem LPS und NINS Dosis
reduzieren, vor der 44. SSW (unreife Leber).
  \item[Rx] 15\,\nicefrac{mg}{kg} (Sgl: 125, KK: 250, Sk: 500, Erw: 1000\,mg
  Supp).
  \item[Cave] ---
  \item[Anm] Antidot: N-Acetylcystein: 140\,\nicefrac{mg}{kg}, dann
  70\,\nicefrac{mg}{kg} alle 8\,h.
\end{labeling}



\subsection{Sonstige}


%%%%%%%%%%%%%%%%%%%%%%%%%%%%%%%%%%%%%%%%%%%%%%%%%%%%%%%%%%%
%%%%%%%%%%%%%%%%%%%%%%%%%%%%%%%%%%%%%%%%%%%%%%%%%%%%%%%%%%%
\subsubsection{4-DMAP (Dimethylaminophenol)}


\begin{labeling}[:]{mm}
  \item[Präp] ---
  \item[Beschreibung] ---
  \item[Ind]  Intox.: Cyanide, Blausäure,
Schwefelwasserstoff.
  \item[KI] Sulfit-Unverträglichkeit (allerg. Asthma), Mischintox mit \ce{CO}
  \item[Rx] initial \VonBis{3}{4}\,\nicefrac{mg}{kg}, anschließend
  \VonBis{100}{500}\,mg Na-thiosulfat.
  \item[Cave] ---.
  \item[Anm] bei Überdosierung (Methämoglobinbildner): Toludinblau: 2\,\nicefrac{mg}{kg}.
\end{labeling}

\section{Häufige Vor- und Dauermedikationen von Patienten}


